\documentclass[11pt]{article}

\usepackage{amsmath,amssymb,amsthm}
\usepackage[margin=20mm]{geometry}

\setlength{\parindent}{0pt}
\setlength{\parskip}{4pt}

\newcommand{\Z}{\mathbb Z}
\newcommand{\Q}{\mathbb Q}
\newcommand{\Hom}{\operatorname{Hom}}
\newtheorem{theorem}{Theorem}
\newtheorem{lemma}[theorem]{Lemma}
\newtheorem{corollary}[theorem]{Corollary}
\theoremstyle{remark}
\newtheorem*{remark}{Remark}

\begin{document}

\begin{center}
{\Large\bfseries The certificate from a symplectic extension}
\end{center}

Let \(q\geq1\), \(C=\Z/q\), and let \(A\) be any finite abelian
group whose exponent divides \(q\).  We identify \(C\) with
\(\frac1q\Z/\Z\subset\Q/\Z\), and put
\(A^\lor=\Hom(A,C)\).  Evaluation gives the canonical
identification \(A=A^{\lor\lor}\).

\begin{lemma}[Polarization]\label{lem:polarization}
If \(E\) is a finite abelian group with an alternating form
\[
             \Omega:E\times E\longrightarrow\Q/\Z ,
\]
then there is a bilinear form \(H:E\times E\to\Q/\Z\) such that
\[
                         H^\top-H=\Omega .                    \tag{1}
\]
\end{lemma}

\begin{proof}
Choose a cyclic decomposition
\[
                         E=\bigoplus_{r=1}^s\langle e_r\rangle .
\]
For \(r<t\), set
\[
 H(e_r,e_t)=0,\qquad H(e_t,e_r)=\Omega(e_r,e_t),\qquad
 H(e_r,e_r)=0,
\]
and extend bilinearly.  This is well defined because the orders of
both \(e_r\) and \(e_t\) annihilate \(\Omega(e_r,e_t)\).  The defining
identity is immediate on the generators.  Notice that no division by
\(2\) occurs.
\end{proof}

\begin{theorem}\label{thm:certificate}
Let \(F\) be a finite-rank free abelian group, let
\(d:F\to F\) be injective with finite cokernel, and let
\[
             F\xrightarrow{\,d\,}F\xrightarrow{\,\pi\,}A
\]
be a complex (\(\pi d=0\)), not necessarily exact.  Let
\(b:F\to A^\lor\), and put
\[
             \varphi=b_\pi,\qquad
             \varphi(x,z)=b(x)\bigl(\pi z\bigr).
\]
Suppose that \(\varphi\) has an integer lift \(\widehat\varphi\) such
that
\[
 \widehat\varphi(x,dy)+\widehat\varphi(y,dx)=0
 \qquad(x,y\in F).                                           \tag{2}
\]
Then there are \(\alpha:F\to A\), a symmetric bilinear form
\(\lambda:A^\lor\times A^\lor\to C\), and an integer lift
\(\widehat\psi\) of \(\psi=b_\alpha\) satisfying
\[
 \partial_d\widehat\psi=\widehat\varphi_d,
 \qquad
 \alpha d=\lambda_*b.                                       \tag{3}
\]
\end{theorem}

\begin{proof}
The map \(d\) becomes an isomorphism over \(\Q\).  Extend
\(\widehat\varphi\) \(\Q\)-bilinearly and define a rational bilinear
form \(\vartheta\) on \(F\otimes\Q\) by
\[
                  \vartheta(dx,z)=\frac{\widehat\varphi(x,z)}q .
                                                                    \tag{4}
\]
Equation (2) says that
\(\vartheta(dx,dy)=-\vartheta(dy,dx)\).  Since \(dF\) spans
\(F\otimes\Q\), the form \(\vartheta\) is alternating.

Form the extension
\[
 \mathcal E=\frac{A^\lor\oplus F}
          {\{(b(x),-dx):x\in F\}},
 \qquad
 i(\chi)=[\chi,0],\quad j(z)=[0,z].
                                                                    \tag{5}
\]
Thus \(jd=ib\) and
\[
                 0\longrightarrow A^\lor
                 \xrightarrow{i}\mathcal E\longrightarrow
                 \operatorname{coker}(d)
                 \longrightarrow0
                                                                    \tag{6}
\]
is exact; in particular \(\mathcal E\) is finite.  On
\(A^\lor\oplus F\), set
\[
\Omega_0\bigl((\chi,z),(\eta,w)\bigr)
 =\chi(\pi w)-\eta(\pi z)+\vartheta(z,w)\pmod{\Z}.             \tag{7}
\]
For every \(x\in F\),
\[
\Omega_0\bigl((b(x),-dx),(\eta,w)\bigr)
 =b(x)(\pi w)+\eta(\pi dx)
  -\frac{\widehat\varphi(x,w)}q=0
 \quad\text{in }\Q/\Z.
\]
The form \(\Omega_0\) is alternating, so the relation subgroup is
also annihilated in the second variable.  Hence (7) descends to an
alternating form \(\Omega\) on \(\mathcal E\).

Apply Lemma~\ref{lem:polarization} and choose
\[
                   H:\mathcal E\times\mathcal E\longrightarrow\Q/\Z,
                   \qquad H^\top-H=\Omega .                  \tag{8}
\]
Define \(\alpha:F\to A=A^{\lor\lor}\) and
\(\lambda:A^\lor\times A^\lor\to C\) by
\[
\boxed{\quad
 \chi(\alpha z)=H(jz,i\chi),\qquad
 \lambda(\chi,\eta)=H(i\eta,i\chi).
\quad}                                                        \tag{9}
\]
These values lie in
\((\Q/\Z)[q]=\frac1q\Z/\Z=C\), because each defining value of \(H\)
has an argument in \(i(A^\lor)\), which is killed by \(q\).  Since
\(\Omega\) vanishes on
\(i(A^\lor)\), equation (8) shows that \(\lambda\) is symmetric.
Moreover, for \(y\in F\) and \(\chi\in A^\lor\),
\[
\begin{aligned}
 \chi(\alpha dy)
 &=H(jdy,i\chi)
  =H(ib(y),i\chi)\\
 &=H(i\chi,ib(y))
  =\lambda\bigl(b(y),\chi\bigr)
  =\chi\bigl(\lambda_*b(y)\bigr).
\end{aligned}
\]
Thus \(\alpha d=\lambda_*b\), which is (P2).

It remains to obtain the prescribed \emph{integer} equality in
(P1).  Choose a rational bilinear lift
\[
        P:F\times F\longrightarrow\Q,\qquad
        P(z,w)\pmod{\Z}=H(jz,jw).
\]
Equations (7)--(8) imply
\[
                     P^\top-P\equiv\vartheta\pmod{\Z}.         \tag{10}
\]
Put \(\delta=P^\top-P-\vartheta\).  This is an integral alternating
form.  Choose an integral bilinear form \(K\) with
\(K^\top-K=-\delta\) (put its entries into one triangle of a matrix).
Replacing \(P\) by \(P+K\), we may therefore assume
\[
                            P^\top-P=\vartheta                \tag{11}
\]
exactly over \(\Q\).

Finally set
\[
                    \widehat\psi(x,z):=qP(z,dx).               \tag{12}
\]
This is integer-valued: modulo \(\Z\),
\[
 P(z,dx)=H(jz,jdx)=H(jz,ib(x))
        =b(x)(\alpha z)\in\tfrac1q\Z/\Z.
\]
The same identity shows that (12) reduces modulo \(q\) to
\(\psi=b_\alpha\).  By (4), (11), and (12),
\[
\begin{aligned}
\widehat\psi(x,dy)-\widehat\psi(y,dx)
 &=q\bigl(P(dy,dx)-P(dx,dy)\bigr)\\
 &=q\,\vartheta(dx,dy)
  =\widehat\varphi(x,dy).
\end{aligned}
\]
This is precisely \(\partial_d\widehat\psi=\widehat\varphi_d\), so
(P1) holds.
\end{proof}

\begin{corollary}[The original matrix hypothesis]\label{cor:matrix}
Let \(R=\Z/q\), let
\[
 D=\operatorname{diag}(d_1,\ldots,d_m),\qquad
 E=\operatorname{diag}(e_1,\ldots,e_n),
 \qquad B,G\in M_{m\times n}(R),
\]
where the \(d_i,e_k\) are positive integers, \(D,E\) denote their
reductions modulo \(q\), and \(d_i\mid d_j\) for \(i<j\).  Put
\(\Gamma=BG^\top\) and assume
\[
 DG=0,\qquad GE=0,\qquad
 \Gamma_{ii}=0,\qquad
 \Gamma_{ji}=-\frac{d_j}{d_i}\Gamma_{ij}\quad(i<j).           \tag{13}
\]
Then there are \(\alpha\in M_{m\times n}(R)\) and
\(\Lambda=\Lambda^\top\in M_n(R)\) such that
\begin{align*}
 (B\alpha^\top)_{ji}-\frac{d_j}{d_i}(B\alpha^\top)_{ij}
     &=\Gamma_{ji} &&(i<j),\\
 D\alpha&=B\Lambda,&
 \alpha E&=0,&
 E\Lambda&=0.
\end{align*}
Thus \textup{(P1)--(P4)} of the original formulation always hold.
Choosing integer representatives gives the integer matrices required
there.
\end{corollary}

\begin{proof}
Use row vectors and put
\[
 A=\{a\in R^{1\times n}:aE=0\}.
\]
The dot product induces a perfect pairing
\[
 \frac{R^{1\times n}}{R^{1\times n}E}\times A\longrightarrow R,
 \qquad(\bar v,a)\longmapsto va^\top;                         \tag{14}
\]
this is the elementary identity
\(\operatorname{ann}(\operatorname{ann}(e_k))=(e_k)\) in
\(\Z/q\), applied coordinatewise, together with equality of the
finite cardinalities.  Hence the first factor in (14) is \(A^\lor\).

Take \(F=\Z^m\), let \(d\) have matrix \(D\), and, writing \(B_i,G_i\)
for rows, define
\[
                 \pi(x_i)=G_i\in A,\qquad
                 b(x_i)=\overline{B_i}\in A^\lor .
\]
The equation \(GE=0\) gives \(\pi(x_i)\in A\), while \(DG=0\) gives
\(\pi d=0\); the map \(b\) is defined by passage to the quotient in
(14).  Moreover \(b_\pi\) has matrix \(BG^\top=\Gamma\).

For \(i<j\), choose an integer \(r_{ij}\) representing
\(\Gamma_{ij}\), and define an integer matrix by
\[
 \widehat\varphi_{ij}=r_{ij},\qquad
 \widehat\varphi_{ji}=-\frac{d_j}{d_i}r_{ij},\qquad
 \widehat\varphi_{ii}=0.                                    \tag{15}
\]
The last two equations in (13) say that (15) is an integer lift of
\(\Gamma\), and
\[
 d_j\widehat\varphi_{ij}+d_i\widehat\varphi_{ji}=0.
\]
Theorem~\ref{thm:certificate} therefore applies.

Write the row \(\alpha_i\) for the element \(\alpha(x_i)\in A\), and
put
\[
 \Lambda_{kl}
   =\lambda(\bar e_k,\bar e_l),
\]
where \(\bar e_k\) are the standard generators of the quotient in
(14).  Then \(\alpha E=0\), \(\Lambda=\Lambda^\top\), and
\(E\Lambda=0\).  Under (14), the element
\(\lambda_*(\bar v)\in A\) is represented by the row \(v\Lambda\).
Consequently \(\alpha d=\lambda_*b\) is exactly
\(D\alpha=B\Lambda\).

Finally \(b_\alpha\) has matrix \(B\alpha^\top\).  On \(x_i,x_j\),
the exact identity furnished by the theorem is
\[
 d_j\widehat\psi_{ij}-d_i\widehat\psi_{ji}
       =d_j\widehat\varphi_{ij}.
\]
After division by \(d_i\), this becomes
\[
 \widehat\psi_{ji}-\frac{d_j}{d_i}\widehat\psi_{ij}
       =-\frac{d_j}{d_i}\widehat\varphi_{ij}
       =\widehat\varphi_{ji}.
\]
Reduction modulo \(q\) is precisely the first displayed equation of
the corollary.
\end{proof}

\begin{remark}
If \(F\xrightarrow dF\xrightarrow\pi A\) is exact, the form
\(\Omega\) is perfect, so \(\mathcal E\) is canonically a finite
symplectic extension of \(A\) by \(A^\lor\).  Exactness and perfection
are not needed for the proof: an adapted Lagrangian complement is also
unnecessary.  The polarization \(H\) is noncanonical, but existence is
all that the certificate requires.
\end{remark}

\begin{remark}
Theorem~\ref{thm:certificate} proves
\texttt{reduced\_conjecture.tex}, while
Corollary~\ref{cor:matrix} proves the original formulation directly.
The argument does not use that \(q\) is a power of \(2\).
\end{remark}

\end{document}
